\begin{textsample}{NEG Dim 4 – global\_north\_2024\_06 – Score 1.00 – global\_north\_2024\_06/200546561.txt}  \label{ex:f4_neg_322}
Anti-hate group ADL slams Wikipedia after site labels it ' unreliable ' source on conflict WASHINGTON -- A leading anti-hate organization that tracks reports of antisemitism is fighting back after a group of anonymous volunteer editors of the online encyclopedia Wikipedia declared the group"generally unreliable"to provide information on the Israel-Palestinian conflict."In a moment when we need to be educating millions of people on how dangerous anti-Jewish hate is, we are being silenced and our research is being marginalized,"Greenblatt said."Just let that sink in."The ADL has tracked anti-Jewish tropes since 1964.
It also researches and documents extremist trends, ideologies and groups across the ideological spectrum.
The global organization has long been considered a bastion for tracking all forms of hate.
News organizations, including USA Today, often cite the ADL ' s research in news articles on antisemitism and extremism.
The decision by a group of Wikipedia editors last week adds the ADL to the list of banned or partially banned sources and puts it in the same company as Russian state @ @ @ @ @ @ @ @ @ @ first reported by the Jewish Telegraph Agency, means the ADL generally should not be cited in articles on the Israel-Palestinian conflict"due to significant evidence that the ADL acts as a pro-Israeli advocacy group and has repeatedly published false and misleading statements as fact,"says a statement posted on a Wikipedia page listing the site ' s frequently used sources.
ADL ' s unreliability extends"to the intersection of the topics of antisemitism and the Israel/Palestine conflict,"the statement said.
Wikipedia ' s review of ADL ' s reliability began in April and involved more than 120 volunteers who participated in discussions about the organization ' s coverage of the Israeli-Palestinian conflict, antisemitism and its hate symbols database, said Maggie Dennis, a vice president at The Wikimedia Foundation, the nonprofit organization that operates the online encyclopedia but does not determine what content is included on the site.
However, the process has attracted attention because of what critics say is a lack of transparency.
The online encyclopedia is volunteer-run and non-profit.
Any reader can see the @ @ @ @ @ @ @ @ @ @ who the volunteer editors are, how they are vetted or if such a process exists at all.
Yet for all its anonymity, the crowd-sourced encyclopedia has over 60 million entries and is often the internet ' s most accessible -- if not always accurate -- record of people, places and events.
In 2023, the English-language version of Wikipedia had about 92 billion views.
In an online discussion forum about ADL, a group of Wikipedia editors suggested the organization is heavily biased in favor of Israel and inaccurately labels legitimate criticism of Israel as antisemitism."The ADL no longer appears to adhere to a serious, mainstream and intellectually cogent definition of antisemitism, but has instead given into the shameless politicisation ( sic ) of the very subject that it was originally esteemed for being reliable on,"wrote an editor who goes by the online name of Iskandar323.
Another editor who uses the online handle Loki said the ADL often acts as a"pro-Israel lobbying organization,"which"can and does @ @ @ @ @ @ @ @ @ @ organizations that disagree with it on this issue."The ADL describes itself as a nonpartisan organization that fights antisemitism, extremism and all forms of hate regardless of ideology or party.
But the group has at times come under attack from both the political left and right for various positions it has taken through the years.
Greenblatt also has faced criticism for suggesting that anti-Israel protesters on college campuses are Iranian"proxies."( In April, the Palestinian militant group Hamas and Iran ' s Supreme Leader Ayatollah Ali Khamenei publicly applauded the growing number of anti-Gaza war encampments that sprung up on American college campuses and became a flashpoint in cities and universities nationwide. ) Greenblatt also has come under fire for comments that critics said equate the Palestinian keffiyeh head scarf to the swastika.
Greenblatt said in an interview with USA Today that"bad actors"with an agenda have intentionally misrepresented his remarks.
His statement about the keffiyeh, for example, came during a television interview in March.
Greenblatt said people @ @ @ @ @ @ @ @ @ @ a swastika on their arm"should not tolerate threats from someone"wearing a keffiyeh on their head.""What I said was that no matter what you ' re wearing, a keffiyeh or anything else, if you make death threats, then you should be accountable for that,"Greenblatt told USA Today."I ' m certainly far from perfect,"he said,"but it is stunning that someone would take comments that were intentionally misrepresented and use that to slander and undermine an entire organization and to suggest that the body of our work is somehow no longer made of integrity."Greenblatt accused Wikipedia of a lack of transparency about its review process and noted that the ADL was never formally notified of the editors ' decision.
The first ADL heard about the ruling was when it was contacted by news organizations, he said.
The Wikipedia decision is problematic for ADL because it raises questions about the organization ' s credibility, said James Loeffler, a professor of modern Jewish @ @ @ @ @ @ @ @ @ @ need good, reliable information about antisemitism, where hate crimes are occurring, and the ADL has long provided critical research,"Loeffler said."The problem is that there ' s one message coming from the leadership and one coming from the research team.
The leadership statements are much less disciplined and raise questions about exactly how rigorous their measurement is and their standards for defining antisemitic acts."If Wikipedia is no longer going to rely on ADL for information about antisemitism and the Israeli-Palestinian conflict,"it ' s a big challenge for having credible information out there for scholars, educators and policymakers as well as the public,"Loeffler said.
Greenblatt said the ADL will not be daunted by Wikipedia ' s decision and will continue its work.

% matched lemmas: 
\end{textsample}
